\chapter{Recurrent neural network e LSTM}
\label{cap:rnn}
\fonti{\rif{2026-10-12}{12/10}, \rif{2026-10-14}{14/10}}{3-deep-learning-methodologies, pp.~34--76; notebook RNN (1-RNN-Pollution-Multivariate, 2-RNN-LSTM-FordA)}

Le RNN estendono la FFNN (cap.~\ref{cap:ffnn}) ai dati \textbf{sequenziali}. Il capitolo va dalla teoria (cella, BPTT,
gradiente, LSTM) alla progettazione e all'implementazione in PyTorch, con i due notebook mostrati a lezione.

% =====================================================================
\section{Perché servono le RNN}
Dato un solo punto (posizione, raggio, colore\ldots) non si può predire la sua \textbf{traiettoria}: manca la storia.
Una FFNN riceve un \textbf{vettore di dimensione fissa} e restituisce un vettore di dimensione fissa; i punti sono trattati
in modo indipendente e i pesi si aggiornano indipendentemente dall'ordine. Se invece do la \textbf{sequenza} delle posizioni
precedenti a una rete che gestisce sequenze, posso indovinare la direzione.

\begin{definizione}[recurrent neural network]
Classe di reti che gestisce input di \textbf{lunghezza variabile} e dati sequenziali $x(t)$. A differenza di una FFNN
mantiene uno \textbf{stato} che \textbf{evolve nel tempo} e descrive l'andamento della serie.
\end{definizione}
Esempio: ``my brother plays \ldots'': sapendo cosa c'è prima, mi aspetto un gioco o uno strumento. In modo
semplificatissimo, è l'idea di predire la parola successiva (alla base di GPT).

\textbf{Applicazioni}: traduzione automatica (una parola dipende dalle precedenti e non sempre corrisponde a una sola parola
nell'altra lingua), previsione di serie temporali (se il valore reale si discosta molto dalla previsione: anomalia, forse un
attacco), classificazione di serie, riconoscimento vocale, \en{part-of-speech tagging} (capire che ``volleyball'' è un nome).

% =====================================================================
\section{FFNN vs RNN in breve}
\begin{itemize}
  \item Una FFNN riceve $x$ e restituisce $y$; il punto successivo viene elaborato indipendentemente.
  \item Una RNN aggiunge un \textbf{self-loop}: riceve $x(t)$ \textbf{e} lo stato nascosto precedente $h(t-1)$, e restituisce
        l'output $y(t)$ e il nuovo stato $h(t)$.
  \item Servono \textbf{nuovi pesi}: oltre a $w$ (importanza delle feature) c'è $w'$ (importanza del passato).
        La complessità aumenta.
  \item Anche le RNN si organizzano in più layer (più layer = rappresentazione più astratta); in pratica 1--2 layer.
\end{itemize}
\begin{definizione}[hidden state]
Rappresentazione interna della serie fino a $t-1$. L'output dipende da $x_t$ \textbf{e} da $h_{t-1}$; il nuovo $h_t$
è aggiornato con l'informazione presente.
\end{definizione}

\begin{center}
\begin{tikzpicture}[every node/.style={font=\small}, cell/.style={draw=polblu, fill=polblu!8, rounded corners=2pt,
  minimum width=1.1cm, minimum height=0.8cm}]
  % forma compatta
  \node[cell] (c) at (0,0) {RNN};
  \node (x) at (0,-1.3) {$x_t$}; \node (y) at (0,1.3) {$\hat y_t$};
  \draw[-{Stealth}] (x) -- (c); \draw[-{Stealth}] (c) -- (y);
  \draw[-{Stealth}] (c.east) to[out=0,in=-90,looseness=1] ++(0.7,0.6) to[out=90,in=0] ++(-0.7,0.5)
        node[above right, font=\scriptsize] {$h_{t-1}$} to[out=180,in=90] ++(-0.7,-0.5) to[out=-90,in=180] (c.west);
  \node at (1.8,0) {$=$};
  % srotolata
  \foreach \i/\t in {0/0,1/1,2/2,3/t} {
    \node[cell] (u\i) at (3+1.8*\i,0) {RNN};
    \node (ux\i) at (3+1.8*\i,-1.3) {$x_{\t}$};
    \node (uy\i) at (3+1.8*\i,1.3) {$\hat y_{\t}$};
    \draw[-{Stealth}] (ux\i) -- (u\i); \draw[-{Stealth}] (u\i) -- (uy\i);
  }
  \draw[-{Stealth}] (u0) -- node[above, font=\scriptsize] {$h_0$} (u1);
  \draw[-{Stealth}] (u1) -- node[above, font=\scriptsize] {$h_1$} (u2);
  \draw[-{Stealth}, dashed] (u2) -- node[above, font=\scriptsize] {$\cdots$} (u3);
\end{tikzpicture}
\end{center}
A sinistra la forma compatta con il self-loop, a destra la versione \textbf{srotolata} (\en{unrolled}): stessa cella, stessi
pesi, ripetuta nel tempo.

% =====================================================================
\section{Varianti input/output}
\begin{center}
\small
\renewcommand{\arraystretch}{1.25}
\begin{tabularx}{\textwidth}{>{\bfseries}l X X}
\toprule
Tipo & Descrizione & Esempi\\
\midrule
One-to-one & un timestamp alla volta, senza passato: è una rete standard & classificazione o regressione su un singolo campione\\
Many-to-one & una sequenza in input, \textbf{un} output; lo stato finale codifica l'intera sequenza & sentiment analysis (``terrifically'' bad/good: serve tutta la frase), sequenza di log per rilevare un attacco in corso\\
One-to-many & un input espanso su più passi di output & image captioning (``un uomo che gioca a baseball''), piano di risposta a partire da un alert\\
Many-to-many & sequenza $\to$ sequenza, lunghezze anche diverse & traduzione, classificazione di ogni frame del traffico in tempo reale\\
\bottomrule
\end{tabularx}
\end{center}
Nel lab ci si concentra sul \textbf{many-to-one}: una serie in input, una label (es.\ malware o no) in output.

% =====================================================================
\section{Architetture: mono e bidirezionale}
\subsection{RNN mono-direzionale}
Elabora i dati in sequenza, dal passato al futuro, aggiornando lo stato con l'input corrente e lo stato precedente.
\begin{itemize}
  \item \textbf{Pro}: adatta alla \textbf{predizione in tempo reale}; relativamente semplice da addestrare (un solo layer, pochi pesi).
  \item \textbf{Contro}: fatica con le \textbf{dipendenze a lungo termine} (dimentica l'inizio di una serie lunga) per
        vanishing/exploding gradient.
  \item \textbf{Training}: \en{Backpropagation Through Time} (BPTT), variante della backprop che tara sia i pesi dell'input sia quelli
        dello stato.
  \item \textbf{Applicazioni}: classificazione supervisionata del traffico (benigno/malevolo), anomaly detection non supervisionata.
\end{itemize}

\subsection{RNN bidirezionale (Bi-RNN)}
Due RNN \textbf{indipendenti} (con \textbf{pesi diversi}, appresi insieme) elaborano la stessa sequenza in direzioni opposte
(avanti e indietro), mantenendo due stati nascosti; gli output vengono poi combinati.
\begin{itemize}
  \item \textbf{Pro}: cattura contesto passato \textbf{e futuro} $\Rightarrow$ rappresentazione migliore (traduzione di una frase completa).
  \item \textbf{Contro}: costo computazionale più alto; \textbf{non adatta al tempo reale} (il futuro non c'è).
  \item Si usa quando si ha \textbf{la sequenza completa}: un'email, quando arriva, è completa $\Rightarrow$ classificazione di
        phishing; riconoscimento di pattern nei log per insider threat.
  \item BPTT applicata separatamente alle due direzioni.
\end{itemize}

% =====================================================================
\section{Struttura interna di una cella RNN}
\[
h_t = F_1\big(W_{hh}h_{t-1} + W_{xh}x_t + b_{hh}\big)
\]
\begin{itemize}
  \item $h_{t-1}$: stato precedente (la ``traiettoria'' fino a $t-1$); $x_t$: input corrente;
  \item $W_{xh}$, $W_{hh}$: matrici dei pesi di input e di stato (ognuna con i propri pesi); $b_{hh}$: bias;
  \item $F_1$: attivazione, tipicamente \textbf{tanh} (non ReLU) per limitare i problemi di gradiente;
  \item per ora l'output della cella $y_t$ coincide con $h_t$.
\end{itemize}

\subsection{Dalla cella alla predizione: la classification head}
$y_t$ (o $h_t$) ha tanti valori quanti la \textbf{hidden size}. Per la predizione finale si aggiunge almeno un layer lineare
(eventualmente dropout/normalizzazione e un'attivazione):
\[
\hat y_t = F_2\big(W_{yh}h_t + b_{yh}\big)
\]
\begin{itemize}
  \item \textbf{Many-to-many}: una head per ogni $x_t$.
  \item \textbf{Many-to-one}: solo l'\textbf{ultima} cella ha la head, perché il suo output ``conosce'' l'intera serie.
\end{itemize}
\begin{esame}[la RNN fa rappresentazione]
La parte ricorrente fa \textbf{feature extraction}: lo stato nascosto alla fine della serie comprime l'andamento della serie in,
ad esempio, 8 numeri: è un \textbf{embedding} della serie. La head (da un layer lineare a una FFNN complessa) usa quella
rappresentazione per classificare. Le due parti si addestrano \textbf{insieme} (ottimizzazione congiunta), non una dopo l'altra.
\end{esame}

\subsection{Dalla FFNN alla RNN, passo per passo}
\begin{enumerate}
  \item Una FFNN con un layer e più output, applicata a un singolo istante $x_t$.
  \item La stessa rete applicata a ogni istante: ogni $\hat y_t$ dipende solo dal suo $x_t$. L'errore è calcolato separatamente per ogni
        rete; $x_0$ e $x_1$ non influenzano $\hat y_2$. \textbf{Nessun concetto di sequenza.}
  \item Si collegano le reti con un link: lo stato $h_t$ viene propagato al passo successivo. Ora
        \[\hat y_t = f(x_t, h_{t-1})\]
        dipende da input presente e memoria passata.
  \item Disegno compatto: la cella con l'arco che torna su se stessa; la catena è la versione \textbf{srotolata} (\en{unrolled}).
\end{enumerate}
\begin{definizione}[relazione di ricorrenza]
\[ h_t = f_W(x_t, h_{t-1}) \]
\textbf{La stessa funzione e lo stesso insieme di pesi} si usano a ogni passo: non sono reti diverse una dopo l'altra, è la stessa
rete ripetuta nel tempo.
\end{definizione}

\subsection{Grafo computazionale nel tempo}
\begin{itemize}
  \item $W_{xh}$ mappa l'input $x_t$ nello stato $h_t$;
  \item $W_{hh}$ definisce come $h_{t-1}$ influenza $h_t$;
  \item $W_{hy}$ mappa lo stato $h_t$ nell'output $y_t$.
\end{itemize}
Tutte queste matrici sono \textbf{condivise tra i passi temporali} e vengono aggiornate con gradienti \textbf{accumulati lungo la
sequenza}.

% =====================================================================
\section{Backpropagation Through Time (BPTT)}
Nella backprop classica (cap.~\ref{cap:reti-neurali}): forward, loss, derivata rispetto a ogni parametro, aggiornamento. Con il tempo:
\begin{enumerate}
  \item si \textbf{sommano le loss} di tutti i passi per ottenere la loss totale della sequenza (in many-to-one c'è una sola predizione,
        ma i pesi vanno comunque ottimizzati lungo tutto il tempo);
  \item l'errore si backpropaga \textbf{sia attraverso gli output sia attraverso gli stati}, dalla fine all'inizio della sequenza;
  \item si \textbf{accumulano i gradienti} per aggiornare i pesi condivisi.
\end{enumerate}
\begin{attenzione}
Per tarare i pesi non basta considerare l'ultimo istante: serve la moltiplicazione con tutti gli stati e gli input precedenti,
indietro nel tempo, anche molto indietro. Per questo è molto più complesso.
\end{attenzione}

\subsection{Il flusso del gradiente}
Il gradiente rispetto a $h_0$ contiene \textbf{molti fattori di $W_{hh}$} e calcoli di gradiente ripetuti. I passi temporali non
sono layer indipendenti (è la stessa rete), ma per il gradiente si comportano come una rete profondissima.
\begin{itemize}
  \item Molti valori $>1$: \textbf{exploding gradient} $\Rightarrow$ i pesi cambiano drasticamente, come con un learning rate enorme;
        apprendimento inaffidabile. Rimedio: \textbf{gradient clipping} (non lasciare crescere il gradiente oltre una soglia).
  \item Molti valori $<1$: \textbf{vanishing gradient} $\Rightarrow$ il passato lontano non conta, la rete non lo ricorda.
        Rimedi: (1) funzione di attivazione, (2) inizializzazione dei pesi, (3) architettura (troppi layer ricorrenti peggiorano).
  \item Nelle RNN si trova spesso \textbf{tanh} invece di ReLU: essendo limitata evita l'esplosione, ma restando $<1$ può far svanire il gradiente.
\end{itemize}
\begin{esame}[exploding vs vanishing nelle RNN]
Il prodotto ripetuto per $W_{hh}$ lungo la sequenza è la causa di entrambi. Exploding $\to$ gradient clipping (e tanh);
vanishing $\to$ attivazione, inizializzazione, architettura, e soprattutto celle con gate (LSTM).
\end{esame}

% =====================================================================
\section{Il problema delle dipendenze a lungo termine}
\begin{esempio}
``The clouds are in the \ldots'': la parola mancante (\emph{sky}) dipende da poche parole vicine. Serve ricorrenza, ma il passato
utile è recente.

``Sono cresciuto in Francia, poi mi sono trasferito in Italia, ho fatto le elementari, giocato a basket, mi sono trasferito a Roma,
ho passato un anno in Australia\ldots'' e dopo un'ora e mezza: ``e parlo fluentemente \ldots''. L'informazione utile (Francia) è
troppo lontana e viene dimenticata. Lo stesso vale per un libro di 1000 pagine in cui conta la prima.
\end{esempio}
\textbf{Mitigazioni}: un'attivazione adatta (ReLU? ovunque? no, altrimenti esplode), una buona inizializzazione, e la soluzione
robusta: una cella più ``intelligente'', la \textbf{gated cell} (LSTM, GRU\ldots; il corso si concentra sull'LSTM). I
\textbf{transformer} (la T di GPT) saranno trattati nella parte NLP.

% =====================================================================
\section{LSTM (Long Short-Term Memory)}
Proposta nel 1997 (Hochreiter e Schmidhuber) per mitigare il vanishing (e l'exploding) gradient.
\begin{itemize}
  \item Mantiene \textbf{due stati}: il \textbf{cell state} $C_t$ (memoria \textbf{a lungo termine}) e l'\textbf{hidden state} $h_t$
        (memoria \textbf{a breve termine}).
  \item Usa dei \textbf{gate} per controllare il flusso: cosa tenere del presente, come aggiornare selettivamente il cell state,
        cosa mandare in output.
  \item Si addestra con BPTT, ma con un flusso di gradiente \textbf{parzialmente ininterrotto} attraverso il cell state.
  \item \textbf{Pro}: gestisce dipendenze a lungo raggio. \textbf{Contro}: molto più costosa di una RNN semplice (più componenti, più
        pesi, più non linearità).
\end{itemize}

\subsection{I quattro passaggi}
\begin{enumerate}
  \item \textbf{Forget gate}: quanto del passato tenere. Prende $h_{t-1}$ e $x_t$, calcola $f_t\in[0,1]$ con una \textbf{sigmoid} e lo
        moltiplica per $C_{t-1}$. $f_t\approx1$: ricorda; $f_t\approx 0$: dimentica (moltiplica per 0 la memoria lunga).
  \item \textbf{Input gate}: quanta nuova informazione aggiungere. $i_t\in[0,1]$ con una sigmoid. $i_t\approx 0$: il passo corrente non
        conta (es.\ un articolo, ``the'', non dice nulla della storia).
  \item \textbf{Update (candidate)}: crea una memoria candidata $\tilde C_t$ con \textbf{tanh}; $i_t$ decide quanta finisce nel cell state.
  \item \textbf{Output gate}: $o_t$ (sigmoid) decide quanto del cell state (passato per tanh) diventa $h_t$, usato per predire $y_t$ o
        passato al passo successivo.
\end{enumerate}
\begin{integrazione}[equazioni standard dell'LSTM]
\begin{align*}
f_t &= \sigma(W_f[h_{t-1}, x_t] + b_f) & i_t &= \sigma(W_i[h_{t-1}, x_t] + b_i)\\
\tilde C_t &= \tanh(W_C[h_{t-1}, x_t] + b_C) & C_t &= f_t \odot C_{t-1} + i_t \odot \tilde C_t\\
o_t &= \sigma(W_o[h_{t-1}, x_t] + b_o) & h_t &= o_t \odot \tanh(C_t)
\end{align*}
$\odot$ = prodotto elemento per elemento. Tre delle attivazioni sono sigmoid: qui servono funzioni \textbf{limitate} perché fanno da
``rubinetti'' tra 0 e 1.
\end{integrazione}
\textbf{Quanto passato ricorda un'LSTM?} Non si sa in anticipo: dipende dal training, dai dati (se il lungo termine serve, la rete lo
terrà) e dalla dimensione di hidden e cell state.

% =====================================================================
\section{Preparare i dati per una RNN}
\label{sec:rnn-dati}
I dati devono essere numeri \textbf{e} sequenze, rappresentati come \textbf{tensore 3D}
\code{(batch, seq\_len, features)}.
\begin{esempio}
\code{[[[1,2,3],[4,5,6]], [[7,8,9],[10,11,12]]]}
\begin{itemize}
  \item ogni vettore (\code{[1,2,3]}) è un istante $x(t)$ e ogni elemento una feature;
  \item ogni matrice (\code{[[1,2,3],[4,5,6]]}) è una serie: 2 istanti descritti da 3 feature;
  \item l'insieme delle matrici è il batch di serie.
\end{itemize}
\end{esempio}
\begin{attenzione}[serie mono-dimensionali]
\code{[1,2,3]} per PyTorch è \textbf{un istante con 3 feature}. Una serie di 3 istanti con una sola feature si scrive
\code{[[1],[2],[3]]}: ogni elemento in un vettore (\code{x.unsqueeze(-1)}, sez.~\ref{sec:rnn-varianti}).
\end{attenzione}

\subsection{Da dati tabellari a sequenze}
Dataset di inquinamento con una riga per ora (35.040 righe, 7 feature + label): ogni riga è un istante, non c'è ancora nessuna serie.
Una funzione \code{create\_sequences} costruisce le serie:
\begin{itemize}
  \item si sceglie la lunghezza (\code{seq\_length = 7}) e si prendono 7 righe consecutive;
  \item la label è quella dell'elemento \textbf{successivo} alla finestra (si predice l'inquinamento all'ora dopo);
  \item si ottengono 35.033 matrici, ognuna con 7 istanti descritti da 7 feature: il tensore 3D per il training.
\end{itemize}
\begin{lstlisting}
def create_sequences(features, labels, seq_length=7):
    X, y = [], []
    for i in range(len(features) - seq_length):
        X.append(features[i:i + seq_length])   # finestra di 7 istanti
        y.append(labels[i + seq_length])       # valore successivo
    return np.array(X), np.array(y)            # X: (N, 7, n_features)
\end{lstlisting}

\subsection{Serie di lunghezza diversa}
In una FFNN un punto con 4 feature \textbf{non} può entrare in un input layer da 5: il quinto neurone non riceverebbe nulla.
Nelle RNN:
\begin{itemize}
  \item il \textbf{numero di feature} per istante deve essere \textbf{sempre lo stesso};
  \item il \textbf{numero di passi temporali può variare} tra serie: un grande vantaggio.
\end{itemize}
Dentro lo stesso \textbf{mini-batch} però le serie devono avere la stessa lunghezza: si usa il \textbf{padding} delle più corte
oppure si \textbf{tronca} tutto a una lunghezza fissa. In generale non serve paddare tutto il dataset: \textbf{solo dentro il batch}
(sez.~\ref{sec:rnn-padding}).

% =====================================================================
\section{Progettare una RNN}
\begin{enumerate}
  \item \textbf{Tipo}: many-to-one, one-to-many, many-to-many? Quanto sono lunghe le sequenze (5 o 5000)? Serve memoria lunga (LSTM)?
        Ha senso leggere anche all'indietro (bidirezionale) o no (mono)?
  \item \textbf{Input}: numero di nodi = numero di \textbf{feature} (non di passi temporali).
  \item \textbf{Hidden size}: quanti neuroni per rappresentare hidden state (e cell state). Più neuroni = più informazione memorizzabile,
        ma più dati per addestrare e più complessità.
  \item \textbf{Hidden layer}: quanti? Più layer = rappresentazione più astratta ma training più difficile. Attivazione tipica: tanh.
  \item \textbf{Padding}: serve? In generale no, dentro il batch sì. Il padding non va considerato come parte reale della serie.
  \item \textbf{Cosa mettere prima dell'output}: un solo layer lineare o una FFNN?
  \item \textbf{Output}: numero di classi, attivazione in base a problema e loss.
  \item \textbf{Ottimizzazione}: mini-batch (shuffle? solo se le serie sono indipendenti), loss (pesi se sbilanciato), ottimizzatore,
        epoch, learning rate; in caso di overfitting, come per le FFNN (cap.~\ref{cap:regolarizzazione}); se forte, ripartire dai dati.
\end{enumerate}

\begin{center}
\small
\renewcommand{\arraystretch}{1.2}
\begin{tabularx}{\textwidth}{>{\bfseries}l X}
\toprule
Iperparametro & Best practice (dalle slide)\\
\midrule
Architettura & mono o bidirezionale, RNN o LSTM in base a cosa si impara dall'ordine all'indietro e a quanto deve essere lunga la memoria\\
\# layer & 1--2 layer ricorrenti, aumentare se serve\\
Hidden size & potenze di 2 (64, 128, 256), non esagerare\\
Attivazione & tanh negli hidden; softmax/sigmoid/lineare in output\\
Inizializzazione & automatica; manuale possibile anche per hidden/cell state\\
Lunghezze diverse & padding o troncamento\\
Batch size & 32--256\\
Shuffle & se le sequenze non sono correlate tra loro\\
Ottimizzatore, LR & Adam (SGD + momentum per la generalizzazione); partire da 0,001 con scheduler\\
Epoch \& early stopping & monitorare la validation loss; il numero di epoch cambia molto con ottimizzatore e LR\\
Regolarizzazione & dropout (0,2--0,5), come layer separato o dentro la RNN; L2 (0,01)\\
Dopo la RNN & valutare se aggiungere layer lineari/non lineari o altre architetture\\
\bottomrule
\end{tabularx}
\end{center}

% =====================================================================
\section{RNN in PyTorch: lunghezze diverse e padding}
\label{sec:rnn-padding}
Una rete è moltiplicazione di matrici. Con tre serie di 2 feature e lunghezze 3, 2 e 3, al terzo passo una matrice ha un
``buco'': non si può moltiplicare per la matrice dei pesi (\rif{2026-10-14}{14/10}).
\begin{itemize}
  \item \textbf{Tagliare} tutto alla lunghezza minima: si butta informazione.
  \item \textbf{Paddare} con valori inventati (es.\ la media): la rete crederebbe che la serie finisca davvero con quei valori.
\end{itemize}
Serve qualcosa di più furbo: padding \textbf{dinamico}, \textbf{solo dentro il mini-batch}, e la RNN deve \textbf{sapere} quali
istanti sono padding.

\subsection{Dataset personalizzato e \code{collate\_fn}}
\begin{lstlisting}
from torch.utils.data import Dataset, DataLoader
from torch.nn.utils.rnn import pad_sequence, pack_padded_sequence

class TimeSeriesDataset(Dataset):
    def __init__(self, series, labels):      # series: lista di tensori (len_i, n_feat)
        self.series, self.labels = series, labels
    def __len__(self):
        return len(self.series)
    def __getitem__(self, idx):               # una serie e la sua label
        return self.series[idx], self.labels[idx]

def collate_fn(batch):
    seqs, labels = zip(*batch)
    lengths = torch.tensor([len(s) for s in seqs])        # lunghezze originali
    padded = pad_sequence(seqs, batch_first=True, padding_value=0.0)  # pad alla max del batch
    return padded, lengths, torch.tensor(labels)

loader = DataLoader(ds, batch_size=64, shuffle=True, collate_fn=collate_fn)
\end{lstlisting}
\begin{itemize}
  \item \code{collate\_fn} riceve i dati del mini-batch, trova la lunghezza massima \textbf{del batch} e padda le serie più corte.
        Restituisce anche le \textbf{lunghezze originali}.
  \item Si passa al \code{DataLoader}, che la usa \textbf{solo se serve}.
  \item Perché basta paddare nel batch? Tutti i passi temporali condividono gli stessi pesi: serve solo che la matrice del batch sia
        completa per moltiplicarla. Batch diversi possono avere lunghezze diverse.
  \item Il padding vale 0, che potrebbe essere un valore reale: per questo la RNN deve sapere dove inizia il padding.
\end{itemize}

\subsection{Il modello}
\begin{lstlisting}
class RNNClassifier(nn.Module):
    def __init__(self, input_size, hidden_size, num_layers, num_classes):
        super().__init__()
        self.rnn = nn.RNN(input_size,             # n. di feature per istante
                          hidden_size,            # dimensione dello stato (embedding)
                          num_layers=num_layers,  # layer ricorrenti impilati
                          batch_first=True)       # input (batch, seq_len, features)
        self.fc = nn.Linear(hidden_size, num_classes)   # classification head
    def forward(self, x, lengths):
        packed = pack_padded_sequence(x, lengths.cpu(), batch_first=True,
                                      enforce_sorted=False)
        out, hidden = self.rnn(packed)    # hidden: (num_layers, batch, hidden_size)
        return self.fc(hidden[-1])        # stato finale dell'ULTIMO layer
\end{lstlisting}
\begin{itemize}
  \item \code{pack\_padded\_sequence}: dice alla RNN quali istanti sono reali e quali padding, così li \textbf{salta}. Il padding serve
        solo alla matematica (matrice completa) e \textbf{non entra} nello stato finale né nella backpropagation: fondamentale, altrimenti
        si addestrerebbe anche sul padding.
  \item \code{hidden[-1]}: lo stato finale dell'\textbf{ultimo} layer, la rappresentazione più astratta della serie.
  \item \code{batch\_first=True}: senza, PyTorch si aspetta \code{(seq\_len, batch, features)}.
\end{itemize}

\subsection{\code{out} vs \code{hidden}}
\begin{center}
\small
\begin{tabularx}{\textwidth}{l X X}
\toprule
& \code{hidden} (\code{h\_n}) & \code{out}\\
\midrule
Cosa contiene & lo stato finale di \textbf{ogni layer} dopo tutta la serie & gli stati dell'\textbf{ultimo layer} a \textbf{ogni passo}\\
Forma & \code{(num\_layers, batch, hidden\_size)} & \code{(batch, seq\_len, hidden\_size)}\\
Uso & task a livello di \textbf{sequenza}: classificazione con head lineare/FFNN (\code{hidden[-1]}) & task a livello di \textbf{passo}: previsione del valore successivo, many-to-many\\
Con padding & con le packed sequence cattura l'ultimo passo \textbf{valido} & l'ultimo passo valido cambia da serie a serie: va selezionato, non preso alla cieca\\
\bottomrule
\end{tabularx}
\end{center}
\begin{esempio}[la scalata]
Sali in montagna registrando battito, pressione e temperatura. \code{hidden} è la rappresentazione dell'intero viaggio (in cima);
\code{out} ti dà la rappresentazione passo per passo, utile se vuoi prevedere il prossimo valore. Nel Lab 2 le sequenze di API call
descrivono un goodware o un malware: serve la rappresentazione \textbf{dell'intera serie} (cap.~\ref{cap:lab2}).
\end{esempio}

\subsection{Varianti: bidirezionale e LSTM}
\label{sec:rnn-varianti}
\begin{lstlisting}
# Bidirezionale
self.rnn = nn.RNN(input_size, hidden_size, num_layers, batch_first=True,
                  bidirectional=True)
self.fc = nn.Linear(hidden_size * 2, num_classes)     # 2 direzioni -> 2x
...
h = torch.cat((hidden[-2], hidden[-1]), dim=1)        # avanti + indietro
logits = self.fc(h)

# LSTM
self.lstm = nn.LSTM(input_size, hidden_size, num_layers, batch_first=True)
...
out, (h_n, c_n) = self.lstm(packed)     # restituisce hidden E cell state
logits = self.fc(h_n[-1])               # per classificare si usa l'hidden state

# Serie univariata: (batch, time) -> (batch, time, 1)
x = x.unsqueeze(-1)                     # [1,2,3] -> [[1],[2],[3]]
\end{lstlisting}

\begin{attenzione}[errori tipici nel codice RNN -- raccolta dei redattori]
\begin{itemize}
  \item Bidirezionale con \code{nn.Linear(hidden\_size, ...)}: l'input della head è \code{hidden\_size * 2}.
  \item \code{out, h = self.lstm(x)}: l'LSTM restituisce \code{(h\_n, c\_n)}; per classificare si usa \code{h\_n[-1]}, non il cell state.
  \item Serie univariata senza \code{unsqueeze(-1)}: \code{[1,2,3]} viene letto come un istante con 3 feature.
  \item Padding senza \code{pack\_padded\_sequence}: la rete si addestra anche sugli zeri di padding.
  \item \code{out[:, -1, :]} con serie paddate: l'ultimo passo può essere padding; con le packed sequence usare \code{hidden[-1]}.
  \item \code{input\_size} pari alla lunghezza della serie invece che al numero di feature per istante.
  \item \code{shuffle=True} su righe che compongono una serie (si mescolano solo serie intere e indipendenti).
\end{itemize}
\end{attenzione}

% =====================================================================
\section{Notebook 1: inquinamento (regressione multivariata)}
\begin{itemize}
  \item Label: livello di inquinamento (numero) $\Rightarrow$ \textbf{regressione}. 7 feature, una riga per ora, train e test già separati.
  \item \textbf{Guardare i dati prima di usarli}: nessun nullo; la direzione del vento è una stringa (SE, NE, cv\ldots) codificata in
        ordine di bussola con un label encoding (semplice, ma forse il one-hot sarebbe meglio).
  \item Plot di serie e distribuzioni: inquinamento fino a 1000 (outlier?) e vento a 600 km/h (\textbf{errore del sensore}?): va deciso
        cosa tenere.
  \item Validation: \textbf{non} con split casuale (romperebbe le serie: il punto 2 in validation tra 1 e 3 in training). Forse meglio
        creare prima le serie e poi dividerle.
  \item Scaling fittato sul training. Nell'esempio si normalizza anche la label: probabilmente non andrebbe fatto.
  \item Serie di lunghezza 7 tutte uguali $\Rightarrow$ basta \code{TensorDataset}. RNN + layer lineare, \textbf{nessuna attivazione finale}
        (regressione), loss \textbf{MSE}. Il training loop è identico a quello della FFNN (cap.~\ref{cap:ffnn}).
  \item Iperparametri del notebook: hidden size 64, output 1, batch 64, Adam con LR 0,001, 200 epoch.
  \item Le loss delle RNN sono \textbf{più rumorose} di quelle delle FFNN; qui non c'è plateau (andava addestrata di più). L'MSE va
        interpretato rispetto alla scala dei valori (nel notebook, su dati scalati: training da circa 0,0056 a 0,0048, validation
        intorno a 0,0055).
\end{itemize}
\begin{verifica}
Nel notebook la validation è ottenuta con \code{df.iloc[train\_size:]}, cioè con uno split \textbf{temporale} (le ultime righe del
training) fatto \textbf{prima} di creare le sequenze, non con uno split casuale. Le osservazioni a lezione sul rischio dello split
casuale restano valide come avvertimento generale.
\end{verifica}

% =====================================================================
\section{Notebook 2: FordA (classificazione di serie)}
\begin{itemize}
  \item Serie \textbf{univariate} con label che indica se il motore funziona; validation = 20\% del training.
  \item Per mostrare il padding si tolgono 20 elementi alle prime 100 serie (480 invece di 500).
  \item Con \code{batch\_size=100}: il primo batch ha tutte serie da 480 (nessun padding), il secondo da 500.
        Con 128: il primo batch contiene serie da 480 e da 500 $\Rightarrow$ padding a 500; i 20 istanti di padding non contano
        per i pesi.
  \item Modelli confrontati: RNN mono, LSTM e Bi-RNN, tutti con hidden size 64, 2 layer, Adam LR 0,001, \code{CrossEntropyLoss}.
  \item Risultato: \textbf{accuracy 51\%} su un dataset bilanciato, con predizione sempre della classe 0 (``tanto vale lanciare una moneta''),
        con RNN, LSTM e Bi-RNN. Non perché le RNN non funzionino: la head è troppo semplice, forse la hidden size è piccola, il
        training è breve. Su CPU è anche molto lento.
\end{itemize}
\begin{integrazione}[leggere la loss bloccata a 0,693]
Nell'output del notebook train e validation loss restano per tutte le 200 epoch intorno a $0{,}693 = \ln 2$: è la cross-entropy di un
classificatore binario che assegna probabilità $0{,}5$ a entrambe le classi. Una loss ferma esattamente lì segnala che il modello
\textbf{non sta imparando nulla}, prima ancora di guardare l'accuracy.
\end{integrazione}
\begin{verifica}
A lezione si è attribuito il risultato anche a un ``training breve'', ma nel notebook le epoch sono 200 (circa 260--390 s su GPU per
modello). Il problema sembra piuttosto di ottimizzazione/architettura (serie lunghe 500 passi, vanishing gradient, head minima).
\end{verifica}

\begin{sintesi}
\begin{itemize}
  \item RNN: stato nascosto $h_t = F_1(W_{hh}h_{t-1} + W_{xh}x_t + b)$ con \textbf{pesi condivisi} nel tempo; tanh tipica; varianti
        one-to-one, many-to-one (focus del corso), one-to-many, many-to-many; mono (tempo reale) o bidirezionale (sequenza completa).
  \item La parte ricorrente produce un \textbf{embedding} della serie; la head classifica; si addestrano insieme.
  \item BPTT: loss sommate nel tempo, gradiente attraverso output e stati; prodotti ripetuti di $W_{hh}$ $\Rightarrow$ exploding
        (gradient clipping) o vanishing gradient.
  \item LSTM: cell state (lungo termine) + hidden state (breve termine); forget, input, candidate (tanh), output gate (sigmoid).
  \item Dati: tensore \code{(batch, seq\_len, features)}; input size = n.\ di feature; lunghezze diverse $\Rightarrow$ padding nel batch
        con \code{collate\_fn} + \code{pack\_padded\_sequence}; classificazione con \code{hidden[-1]}, previsione passo-passo con \code{out}.
  \item Notebook: pollution (regressione, MSE, nessuna attivazione finale); FordA (51\%: loss ferma a $\ln 2$, il modello non impara).
\end{itemize}
\end{sintesi}
